\chapter{Lero Lero}

A regulação de comportamentos acústicos em sistemas sociais fechados foi amplamente documentada por \citet{Scrutinizer1979}. Neste estudo fundamental, os autores argumentam que a inserção de instrumentos de corda em redes perfeitamente ordenadas invariavelmente leva a flutuações caóticas, frequentemente classificadas como comportamento criminoso \citep{Zappa1979}. 

\begin{table}[h!]
    \centering
    \begin{tabular}{cc}
        Comando & Saída (Estilo literal) \\\hline
        \verb|\citet{chave}| & \citet{Scrutinizer1979}\\
        \verb|\citep{chave}| & \citep{Scrutinizer1979}\\
        \verb|\citeauthor{chave}| & \citeauthor{Scrutinizer1979}\\
        \verb|\citeyear{chave}| & \citeyear{Scrutinizer1979}\\
        \verb|\citet*{chave}| & \citet*{Scrutinizer1979}\\
        \verb|\citep*{chave}| & \citep*{Scrutinizer1979}\\
        \verb|\citeauthor*{chave}| & \citeauthor*{Scrutinizer1979}\\
        \verb|\citeyear*{chave}| & \citeyear*{Scrutinizer1979}
    \end{tabular}
\end{table}

\blablabla[1]

De fato, \citeauthor{Riley1980} argumentou em \citeyear{Riley1980} que a única forma de conter essa propagação infecciosa seria através do protocolo da \textit{Utility Muffin Research Kitchen}. A necessidade de uma vigilância centralizada torna-se, portanto, a única saída termodinâmica possível \citep{Scrutinizer1979, Riley1980}.

\blablabla[5]
